\documentclass[12pt]{article}
\usepackage[utf8]{inputenc}
\usepackage[T1]{fontenc}
\usepackage[margin=1in]{geometry}
\usepackage{setspace}
\usepackage{xcolor}
\usepackage[colorlinks=true, urlcolor=blue!60!black, linkcolor=blue!60!black]{hyperref}
\setlength{\parindent}{2em}
\setlength{\parskip}{1.25em}
\renewcommand{\baselinestretch}{1.0}
\begin{document}
\begin{center}
{\LARGE \textit{United States v. Texas}, 599 U.S. 670 (2023)\par}
\vspace{0.8em}
{\large Case Brief --- Four-Page Edition\par}
\end{center}
\vspace{1em}

\section{Memory Jogger}

States lack Article III standing to sue the federal Executive Branch to compel more arrests or prosecutions of third-party noncitizens, because indirect fiscal costs from non-enforcement are not a traditionally judicially cognizable injury (\href{https://www.law.cornell.edu/supremecourt/text/22-58}{Cornell LII}).

\section{Detailed Case Facts}

With more than 11 million removable noncitizens and limited resources, administrations have long set immigration enforcement priorities. In September 2021, DHS Secretary Mayorkas issued Guidelines prioritizing noncitizens who threaten national security, public safety, or border security, and directing agents to weigh aggravating and mitigating factors rather than rely solely on a criminal conviction. Texas and Louisiana sued under the APA, claiming the Guidelines violated 8 U.S.C. \S~1226(c), which says the government ``shall take into custody'' certain criminal noncitizens, and \S~1231(a)(2), which says it ``shall detain'' noncitizens during the removal period. The States alleged costs from continued criminal supervision, recidivism, education, and healthcare.

\section{Procedural History}

After a bench trial, Judge Tipton in the Southern District of Texas found standing, held the Guidelines unlawful, and vacated them nationwide. The Fifth Circuit denied a stay. The Supreme Court denied a stay 5--4 but granted certiorari before judgment on standing, the merits, and whether \S~1252(f)(1) bars vacatur, and it reversed on standing alone.

\section{Judicial Votes}

Reversed, \textbf{8--1} (decided June 23, 2023).
\begin{itemize}
\item \textbf{Majority:} Justice Kavanaugh, joined by Chief Justice Roberts and Justices Sotomayor, Kagan, and Jackson.
\item \textbf{Concurring in the judgment:} Justice Gorsuch, joined by Justices Thomas and Barrett; Justice Barrett, joined by Justice Gorsuch.
\item \textbf{Dissent:} Justice Alito.
\end{itemize}

\section{Holding}

A plaintiff, including a State, lacks a judicially cognizable interest in challenging the Executive's arrest or prosecution policies when it is neither prosecuted nor threatened with prosecution, even if under-enforcement imposes indirect costs. The Court did not decide whether the Guidelines comply with the statutes or whether universal APA vacatur is permitted. It listed five situations that ``might'' differ: selective prosecution, congressionally created injuries and remedies, complete abdication, non-enforcement combined with legal benefits, and continued detention.

\section{Analysis of Opinions}

\textbf{Majority (Kavanaugh).} Monetary costs are an injury in fact, but under \textit{Raines} and \textit{TransUnion} an injury must also be legally and judicially cognizable. No history supports orders compelling more arrests, and \textit{Linda R. S. v. Richard D.} holds that a citizen not prosecuted or threatened cannot contest prosecution policy. Three reasons support the rule: non-enforcement does not coerce the plaintiff; such suits intrude on Article II enforcement discretion, especially in immigration; and courts lack standards to second-guess resource trade-offs (\textit{Heckler v. Chaney}). \textit{Massachusetts v. EPA} involved a rulemaking petition. Political checks remain.

\textbf{Gorsuch (concurring in the judgment).} The defect is redressability: \S~1252(f)(1) bars injunctions, and vacating the Guidelines would leave officials the same discretion. He would shelve special solicitude for States and doubted that \S~706(2) authorizes universal vacatur.

\textbf{Barrett (concurring in the judgment).} The majority overread \textit{Linda R. S.}, a redressability case, and wrongly fused merits and statutory-review doctrines into Article III.

\textbf{Alito (dissenting).} Texas met the \textit{Lujan} test on unchallenged findings; Congress used ``shall'' to remove discretion; and shielding executive disobedience revives the suspending and dispensing powers repudiated in 1689.

\section{Examples: Future Applications}

\begin{enumerate}
\item \textbf{Same-Side (No Standing): Marijuana Prosecutions.} A state challenges DOJ's deprioritizing of simple-possession prosecutions, citing policing costs.
\item \textbf{Same-Side (No Standing): ATF Priorities.} Police unions and counties challenge guidance favoring violent-crime investigations over paperwork violations under a ``shall investigate'' statute.
\item \textbf{Opposite-Side (Standing): Benefits Program.} A DAPA-like program confers lawful presence and work authorization, implicating more than arrest discretion.
\item \textbf{Opposite-Side (Standing): Rulemaking Petition.} A state challenges EPA's denial of a rulemaking petition, governed by \textit{Massachusetts v. EPA}.
\item \textbf{Fence-Sitter: Citizen-Suit Provision.} Congress lets border states sue to compel custody under \S~1226(c). The second caveat suggests standing, but footnote 4 warns of Article II problems.
\end{enumerate}

\section{Critique}

\textbf{Originalist views.} Admirers of Gorsuch praise his Article III rigor and his textual challenge to universal vacatur; Alito's allies and the \textit{Louisiana Law Review} fault the majority for nullifying ``shall'' mandates and giving agencies carte blanche.

\textbf{Progressive and doctrinal views.} Many welcomed the curb on state-attorney-general suits in single-judge districts, but the \textit{Harvard Law Review} criticized a selective use of history that deepens standing law's fragmentation, and scholars echo Barrett's worry about constitutionalizing \textit{Heckler}. The case pits executive flexibility and restraint against legislative supremacy and federalism.

\section{Key Quotations}

\begin{quote}
``The States essentially want the Federal Judiciary to order the Executive Branch to alter its arrest policy so as to make more arrests. But this Court has long held `that a citizen lacks standing to contest the policies of the prosecuting authority when he himself is neither prosecuted nor threatened with prosecution.' \textit{Linda R. S. v. Richard D.}, 410 U. S. 614, 619 (1973). Consistent with that fundamental Article III principle, we conclude that the States lack Article III standing to bring this suit.''
\end{quote}
\noindent --- Justice Kavanaugh, for the Court.

\begin{quote}
``It bears emphasis that the question of whether the federal courts have jurisdiction under Article III is distinct from the question of whether the Executive Branch is complying with the relevant statutes---here, \S 1226(c) and \S 1231(a)(2). In other words, the question of reviewability is different from the question of legality. We take no position on whether the Executive Branch here is complying with its legal obligations under \S 1226(c) and \S 1231(a)(2). We hold only that the federal courts are not the proper forum to resolve this dispute.''
\end{quote}
\noindent --- Justice Kavanaugh, for the Court.

\begin{quote}
``The States lack standing because federal courts do not have authority to redress their injuries. Section 1252(f)(1) denies the States any coercive relief. A vacatur order under \S 706(2) supplies them no effectual relief. And such an order itself may not even be legally permissible. ... The Constitution affords federal courts considerable power, but it does not establish `government by lawsuit.' ''
\end{quote}
\noindent --- Justice Gorsuch, concurring in the judgment.

\end{document}
